% !PRINTED 140
\weberpage{140}
%<<CONT
that a simple permutation group of degree 60 on five figures actually exists, namely the alternating group.

But one can also verify the fact by exhibiting a divisor of degree 12 of $P$ (\S.\,28, 3.). One obtains this divisor of degree 12, e.g., from the fact that the group of degree 4:
\[
Q = 1, (1, 2) (3, 4), (1, 2) (5, 6), (3, 4) (5, 6)
\]
is permutable with the permutations of degree 3:
\[
c = (1, 3, 5) (2, 4, 6)
\]
that is, satisfies the condition $c^{-1} Q c = Q$, and one thus then obtains a group of degree 12:
\[
Q, \; Q c, \; Q c^2.
\]
The simple group of degree 60 is also called the \gesperrt{icosahedral group}, for a reason that we shall come to know later.

The same consideration also shows that a simple group of degree 90 does not exist. For this group would have to contain an element of degree 5 and an element of degree 3. Since $90 = 5 . 18$ and 18 has no divisor other than 6 that is congruent to 1 modulo 5, it follows that, if we choose for $Q$ a group of degree 5, $R$ must be of index 6, and the group $P$ would therefore be representable as a transitive permutation group of six figures. Just as above, one concludes that it would have to contain the two elements $a, b$, which, as we have seen, leads to the icosahedral group, which cannot be a divisor of a group of degree 90.

\begin{center}
\S.\,35.

\gesperrt{Groups of degree $pq$.}
\end{center}

It is now of further interest to investigate how, for a given degree number, the groups are constituted. For this the theorems derived in the preceding sections are the foundations. To be sure, we are as yet in a position to survey the existing groups completely only for the simpler degree numbers. A work by \gesperrt{Hölder}\footnote{\gesperrt{Hölder}, Die Gruppen der Ordnungen $p^3, pq^2, pqr, p^4$. Mathematische Annalen, Bd. 43.} goes furthest in this direction. We consider here only the simple case,
%CONT>>
